%% kapa: exam=1 round=22 year=2011 period=1 form=A subject=civil_law subject_ko=민법 seq=10 number=10
%% source: Q-Net 감정평가사 기출문제 (cst00309) artlSeq=5036790 file=extracted/1008/5036790/2011년 제22회 감정평가사 1차시험 기출문제/1교시.hwp
%% conversion: kordoc 4.24.51 -> tools/gen1.py (auto); review=proofread-sonnet; flags=none
\begin{question}{10}
소멸시효의 중단사유로서의 승인에 관한 설명으로 옳지 않은 것은? (다툼이 있으면 판례에 의함)
\begin{choices}
  \item 승인은 소멸시효의 진행이 개시된 이후에만 가능하고 그 이전에 승인을 하더라도 시효가 중단되지는 않는다.
  \item 변제기한의 유예를 요청하는 것은 채무의 승인에 해당한다.
  \item 시효완성 전에 한 채무의 일부변제는, 그 금액에 관하여 다툼이 없는 한, 채무승인에 해당한다.
  \item 시효중단의 효력이 있는 승인에는 상대방의 권리에 관한 처분의 능력이나 권한을 요한다.
  \item 채무승인이 있었다는 사실은 이를 주장하는 채권자 측에서 증명하여야 한다.
\end{choices}
\end{question}
\answer{4}
